\section{Leakage-Free Temporal Topology Engine}\label{sec:engine}

A transaction does not reveal fraud in isolation; the same transfer can be
ordinary in one account context and the closing leg of a laundering cycle in
another. Graph-based fraud detection formalizes this by treating the ledger as
a network of who-paid-whom and searching for structural motifs — cycles,
fan-in collection points, rapid pass-through chains — that a row-wise view
cannot see~\cite{motie2024gnn,yuan2025survey,salih2021survey}. The difficulty
is computing such motifs \emph{as they would have been visible at the time},
without letting the completed, labeled dataset's hindsight leak into a
feature that is supposed to represent an auditor's knowledge at transaction
time. This section describes the streaming component of the Self-Auditing
Ledger that does this — the Temporal Topology Engine — and states, precisely
and with the project's own evidence, which of its thirteen output features
actually earn their keep.

The engine operates on the canonical edge list produced by ETL: one row per
transaction, with a source account, a destination account, a timestamp, an
amount, and (for training data) a fraud label. It is deliberately \emph{not}
a graph neural network and does not touch Neo4j; it is a single linear pass
over the sorted edge stream, maintained in pure Python/NumPy, that emits one
feature vector per transaction. Two data structures carry all state across
that pass: a trailing-window directed multigraph, \texttt{WindowGraph}, and a
per-account rolling amount memory, \texttt{LapMemory}. Both are described
below, followed by the thirteen features they jointly produce, the
structure/amount split used to keep the later ablation honest, and the
feature-importance result — including a null result the project does not
soften.

\subsection{The Leakage Invariant: Read, Then Insert}\label{sec:engine-invariant}

CLAUDE.md's rule 4 states the requirement in one sentence: a transaction's
features may depend only on transactions with \texttt{timestamp <=} its own.
The engine satisfies this not as a coding convention that a future edit could
quietly violate, but as a structural property of a single loop body. Edges
are pre-sorted by \texttt{(timestamp, transaction\_id)} — the ETL writes them
in this order with a stable sort, and the feature-assembly stage re-sorts
defensively if that invariant is ever violated upstream. The engine then
iterates \texttt{for i in range(n)} and, for transaction $i$, executes
exactly two kinds of statement in exactly one order:

\begin{enumerate}
  \item \textbf{Read.} Query \texttt{WindowGraph} for $i$'s fan-in, fan-out,
  destination in-degree, source out-degree, and (if $\mathrm{src}_i \neq
  \mathrm{dst}_i$) run the bounded cycle search; query \texttt{LapMemory} for
  the recurrence of $i$'s amount against each endpoint's rolling history.
  \item \textbf{Insert.} Only after every read above returns does the engine
  call \texttt{WindowGraph.add(t, u, v, amount)} and
  \texttt{LapMemory.record(\ldots)} for transaction $i$ itself.
\end{enumerate}

Because \texttt{add} for edge $j$ is called only during iteration $j$, and
iterations execute strictly in order $0, 1, \ldots, n-1$, the graph state
visible to iteration $i$'s reads is, by induction, exactly $\{j : j < i\}$
minus whatever \texttt{evict} has aged out — never a superset, and in
particular never anything from an iteration that has not yet run. No branch
of the loop calls \texttt{add} or \texttt{record} ahead of the paired reads
for the same $i$, and \texttt{evict} only removes edges whose timestamp falls
below \texttt{now - window}; it has no code path that admits an edge, only
ones that retire old ones. This is what we mean by \emph{leakage-free by
construction}: the claim does not rest on trusting that every caller remembers
to query before mutating, or on empirically observing no leak on the datasets
at hand — it follows from the fixed order of two statements inside one loop,
for any input that satisfies the pre-sort precondition. \texttt{LapMemory}
enforces the identical discipline: \texttt{recurrence(acct, amount)} is
always read before the same call site's \texttt{record(acct, amount,
counterparty)}. For SAP Würzburg, whose synthetic per-run time axis (encoded
as $\mathrm{run\_index}\times 10^{6} + \text{seconds-since-midnight}$) is not
a real cross-run chronology, both structures are additionally discarded and
rebuilt at every \texttt{run\_id} boundary, so a document in run 214 can never
borrow structural context assembled from a different, unrelated ledger run.

This constructive guarantee is checked, not merely asserted: an automated
leakage test (\texttt{topology.leakage\_test}) re-derives the same feature
table twice on real, built data — once truncated at the halfway point, and
once with every transaction \emph{after} the halfway point corrupted (shuffled
endpoints, amounts overwritten with an out-of-range sentinel) while
timestamps and row order are preserved — and asserts the first half's
features are byte-identical in both cases. This is deliberately redundant
with the constructive argument above: the constructive argument is the proof
that leakage is impossible for any input; the test is a regression guard that
would fail loudly if a future refactor moved a read after an insert. It has
run green throughout the build.

\subsection{The Trailing-Window Graph}\label{sec:engine-windowgraph}

\texttt{WindowGraph} holds a directed multigraph restricted to a trailing
time window $[\mathrm{now} - w, \mathrm{now}]$, where $w$ is a per-dataset
constant (7 days for IBM AML; 7 steps for BankSim, whose control role
requires \texttt{in\_cycle} to be identically zero everywhere since its
graph is customer$\to$merchant bipartite; 24 steps $=1$ day for PaySim; and
unbounded \emph{within} a run for SAP Würzburg, whose runs are short enough
that no window truncation is needed). Internally it keeps adjacency as
per-node dictionaries of neighbor multiplicities, so fan-in/fan-out (count of
\emph{distinct} counterparties) and weighted in-/out-degree (count
\emph{with} multiplicity) are both $O(1)$ lookups against the current
window state, and a deque of \texttt{(timestamp, u, v)} triples in arrival
order lets \texttt{evict(now)} pop expired edges from the front in amortized
$O(1)$ per evicted edge. The one non-trivial query is the cycle search: does
inserting $u \to v$ close a directed loop $v \rightsquigarrow u$ already
present in the window? This is answered by a breadth-first search from $v$
along existing out-edges, capped at \texttt{max\_cycle\_len}$-1$ hops (default
6) and a visited-node \texttt{search\_budget} (default 4{,}000); either bound
being hit returns "no cycle" rather than risk a false positive or an
unbounded cost on a dense hub account, so the detector is a conservative
under-approximation of the true cycle set, never an over-approximation. Each
directed pair also remembers the \texttt{(amount, timestamp)} of its most
recent in-window edge, which the cycle search folds together to report
amount conservation and timing tightness around the closed loop without a
second graph pass.

\subsection{Per-Account Lapping Memory}\label{sec:engine-lapmemory}

Lapping — using a new, similarly sized inflow to cover an earlier shortfall,
"robbing Peter to pay Paul" — is not naturally a graph-window feature: its
tell-tale recurrence can be separated by weeks or months, far longer than any
window size that keeps the cycle search cheap. \texttt{LapMemory} instead
keeps a fixed-length, \emph{count}-based rolling history of the last $k=25$
$(\text{amount}, \text{counterparty})$ pairs an account was involved in, in a
given role (incoming or outgoing) — two independent instances are kept per
account, one for each role. Because eviction is by count, not by elapsed
time, this memory spans exactly as much calendar time as the account's own
activity took to accumulate 25 transactions, so a dormant account that
reappears after a long gap is still compared against its own last known
pattern rather than an empty, window-truncated one. \texttt{recurrence(acct,
amount)} scans that account's stored history for entries within a relative
tolerance $\mathrm{tol}=1\%$ of the query amount and returns both the count of
matches and the number of \emph{distinct} counterparties among them. The
second number is what makes this a lapping detector rather than a
recurring-bill detector: one counterparty repeating the same payment every
month is ordinary billing, but several \emph{different} counterparties each
supplying a near-identical amount to the same account is the diversification
pattern lapping schemes exhibit. Reporting raw recurrence count alone would
fall into the same triviality trap as raw cycle membership (\S
\ref{sec:engine-finding}) — it fires constantly on legitimate recurring
payments — so counterparty diversity is carried as a co-equal, separate
feature rather than folded into a single scalar.

\subsection{The Thirteen Topology Features}\label{sec:engine-features}

Every transaction receives the same thirteen columns, grouped below by the
structure they measure; all are computed from the past-only graph/memory
state exactly as described above.

\textit{Cycle detectors (5).} \texttt{in\_cycle} is $1$ if the transaction
closes some directed loop within the window (a same-account self-payment
counts as a degenerate length-1 loop and is flagged directly, without a
search); \texttt{cycle\_length} is the edge count of the \emph{shortest} such
loop found, $0$ if none; \texttt{in\_cycle\_ge3} restricts this to loops of
length $\geq 3$, excluding both self-loops and simple reciprocal pairs
($A\to B\to A$), which are common and largely benign; \texttt{cycle\_amount\_ratio}
is $\min/\max$ of the amounts around the closed loop (including the closing
edge itself) — near $1$ means the money round-tripped with its value intact,
the discriminating signature of a wash trade rather than, say, a sale
followed by a small partial refund; \texttt{cycle\_time\_span} is the closing
transaction's timestamp minus the earliest edge on the loop — small values
mean the loop closed quickly, the "tightness" that separates a fast
laundering round-trip from an incidental, slow-forming one. Both are
\texttt{NaN} when no cycle closes, left as missing for the tree model rather
than imputed to zero.

\textit{Degree/fan detectors (4).} \texttt{fan\_in} and \texttt{fan\_out}
count, respectively, the number of \emph{distinct} accounts that have sent to
the destination account and the number of \emph{distinct} accounts the
source account has paid, within the window; \texttt{dest\_in\_degree} and
\texttt{src\_out\_degree} count the corresponding edges \emph{with}
multiplicity. The distinct/multiplicity pairing lets the model separate "many
different counterparties" (structuring, collection points) from "one
counterparty transacting repeatedly" (high degree, low fan) — the same
diversity-versus-repetition distinction the lapping features apply to
amounts.

\textit{Lapping detectors (4).} \texttt{lap\_in\_recur} and
\texttt{lap\_in\_payers} are \texttt{LapMemory.recurrence} evaluated against
the destination account's \emph{incoming}-amount history; \texttt{lap\_out\_recur}
and \texttt{lap\_out\_payees} the same against the source account's
\emph{outgoing}-amount history (\S\ref{sec:engine-lapmemory}).

\subsection{Structure vs.\ Amount: Why the Split Matters}\label{sec:engine-split}

Of the thirteen features, five are computed \emph{from transaction amounts}
— \texttt{cycle\_amount\_ratio} and all four lapping features — while the
remaining eight (\texttt{in\_cycle}, \texttt{cycle\_length},
\texttt{in\_cycle\_ge3}, \texttt{cycle\_time\_span}, \texttt{fan\_in},
\texttt{fan\_out}, \texttt{dest\_in\_degree}, \texttt{src\_out\_degree}) are
pure graph shape and never look at an amount value.
\texttt{cycle\_time\_span} is time-derived, not amount-derived, and is
counted with the structure group despite superficially "sounding" adjacent
to the amount features. The engine exposes this split (\texttt{topology\_structure\_cols}
vs.\ \texttt{topology\_amount\_cols} in the feature manifest) because the
ablation study's central probe — does topology help \emph{independent of the
amount signal}? — requires dropping every ordinary amount-derived tabular
feature (\texttt{log\_amount}, the expanding amount z-score, etc.)
\emph{and} every amount-derived topology feature together. Running an
"amount-blind" arm that strips the ordinary amount columns but leaves
\texttt{lap\_*} and \texttt{cycle\_amount\_ratio} in place would silently
re-admit the amount signal through the back door and overstate what pure
graph shape contributes. The structure-only arms in Table~\ref{tab:topology-importance}'s
source data (\texttt{no\_amount\_plus\_structure} in \texttt{paper/data/ablation.json})
are the clean version of this test; their interpretation is developed fully
in the ablation study.

\subsection{Feature Importance, and a Result We Do Not Bury}\label{sec:engine-finding}

Table~\ref{tab:topology-importance} reports each topology feature's XGBoost
gain importance~\cite{chen2016xgboost}, normalized within the full
topology-augmented model (\texttt{topology\_full} arm) for each of the four
datasets carrying a complete ablation run. Because these thirteen columns
sit alongside 15--21 ordinary tabular features in that model (BankSim: 21
ordinary + 13 topology = 34; IBM AML: 8 + 13 = 21; SAP Würzburg: 20 + 13 =
33; synth\_erp: 15 + 13 = 28), the rows below do not sum to $1$ within a
dataset — the remaining gain mass sits with the ordinary features, not shown
here. SAP Würzburg's split sizes in this table (120{,}411 / 40{,}137 /
40{,}139, partition-aware by \texttt{run\_id}) come from
\texttt{ablation.json} specifically; the paper's headline champion-model
table elsewhere draws on a very slightly different SAP split
(120{,}341 / 40{,}114 / 40{,}116) recorded in \texttt{final\_champions.json} — the
two pipelines' SAP preprocessing diverged by 70 rows in the training split
and 23 rows in each of validation and test, and
the two are never mixed within one reported number.

\begin{table}[t]
\centering
\caption{Topology feature gain importance (fraction of total gain in the
topology-augmented model), by dataset. Source:
\texttt{paper/data/ablation.json}, \texttt{topology\_importance}, single
seed, no confidence intervals.}
\label{tab:topology-importance}
\resizebox{\columnwidth}{!}{%
\begin{tabular}{lrrrr}
\toprule
\textbf{Feature} & \textbf{BankSim} & \textbf{IBM AML} & \textbf{SAP Würzburg} & \textbf{synth\_erp} \\
\midrule
\texttt{in\_cycle}           & 0.0000 & 0.0000 & 0.0000 & 0.0040 \\
\texttt{cycle\_length}       & 0.0000 & 0.0009 & 0.0000 & 0.0029 \\
\texttt{in\_cycle\_ge3}      & 0.0000 & 0.0000 & 0.0000 & 0.0044 \\
\texttt{cycle\_amount\_ratio}& 0.0000 & 0.0008 & 0.0000 & 0.0090 \\
\texttt{cycle\_time\_span}   & 0.0000 & 0.0001 & 0.0000 & 0.0072 \\
\midrule
\texttt{fan\_in}             & 0.0908 & 0.0003 & 0.1791 & 0.0241 \\
\texttt{fan\_out}            & 0.0062 & 0.0010 & 0.0136 & 0.0189 \\
\texttt{dest\_in\_degree}    & 0.5199 & 0.0003 & 0.0682 & 0.0230 \\
\texttt{src\_out\_degree}    & 0.0062 & 0.0007 & 0.0089 & 0.0786 \\
\midrule
\texttt{lap\_in\_recur}      & 0.0052 & 0.0006 & 0.0319 & 0.0130 \\
\texttt{lap\_in\_payers}     & 0.0048 & 0.0006 & 0.0199 & 0.0103 \\
\texttt{lap\_out\_recur}     & 0.0066 & 0.5397 & 0.0416 & 0.0041 \\
\texttt{lap\_out\_payees}    & 0.0079 & 0.4200 & 0.0000 & 0.0036 \\
\bottomrule
\end{tabular}%
}
\end{table}

%% FIGURE TODO: grouped bar chart, one panel per dataset (BankSim, IBM AML,
%% SAP Würzburg, synth_erp), bars = gain importance for the 13 topology
%% features colour-coded by group (cycle / degree-fan / lapping), y-axis
%% shared across panels to make visually obvious that the 5 cycle-feature
%% bars are flat at zero in 3 of the 4 panels and barely visible in the 4th.
\begin{figure}[t]
\centering
% \includegraphics[width=\columnwidth]{figures/topology_importance_by_group.png}
\caption{Gain importance of the thirteen topology features, grouped by
detector family and faceted by dataset. Cycle-detector bars (leftmost group
in each panel) are at or indistinguishable from zero in three of four
datasets.}
\label{fig:cycle-null}
\end{figure}

The pattern the table makes plain is not subtle, and we state it without
qualification: the five cycle features — \texttt{in\_cycle},
\texttt{cycle\_length}, \texttt{in\_cycle\_ge3}, \texttt{cycle\_amount\_ratio},
\texttt{cycle\_time\_span} — carry essentially zero gain importance on three
of the four datasets. On BankSim and SAP Würzburg every one of the five is
exactly $0.0000$; on IBM AML the largest is \texttt{cycle\_length} at
$0.0009$, well over two orders of magnitude below that dataset's dominant features
(\texttt{lap\_out\_recur} $0.5397$, \texttt{lap\_out\_payees} $0.4200$). Only
on \texttt{synth\_erp} — our own generated dataset, used for
training/stress-testing and never as headline proof — do cycle features
carry non-trivial weight, and even there the largest single one
(\texttt{cycle\_amount\_ratio}, $0.0090$) is under one percent of total gain.
This is a pointed result for a project whose title is built around a
"shadow ledger" of detected cycles: the cycle detector is not inert by
construction (\S\ref{sec:engine-windowgraph} shows it is a real, bounded BFS
that fires and populates non-\texttt{NaN} values whenever a loop closes,
and it is validated separately against IBM AML's labeled ground-truth cycles
in the ablation study), yet the boosted-tree model consistently prefers
other signals to explain fraud once they are available. IBM AML is the most
striking case, because it is precisely the dataset engineered with labeled
cycle-shaped fraud~\cite{altman2023realistic} — and it is exactly where
cycle-feature gain is most thoroughly crowded out: \texttt{lap\_out\_recur}
and \texttt{lap\_out\_payees} alone account for gain fraction $0.9597$ out of
the full 21-feature model's normalized total of $1.0$ — the two amount-derived
lapping features claim $96\%$ of \emph{every} split the model makes, ordinary
tabular features included, leaving essentially nothing for the five cycle
columns.
Our working hypothesis, developed fully where the predictive consequences of
this pattern are analyzed, is two-part: true cycles are rare enough in every
dataset we could build (a handful of closures against hundreds of thousands
of transactions) that a gradient-boosted tree finds little marginal
information in them once even coarser degree/fan and amount-recurrence
signals are on the table, and where a cycle transaction and an
amount-lapping signature co-occur, the amount/lapping features are simply
more discriminative per split — so the tree spends its gain budget there
first and has little left to reward on the rarer cycle columns. We do not
claim this hypothesis is proven by Table~\ref{tab:topology-importance} alone;
it is a description of what the importances show, offered as the
starting point the ablation study's marginal-lift and amount-blind results
either support or complicate.
